\documentclass[12pt,letterpaper]{article}
\usepackage[utf8]{inputenc}
\usepackage[spanish]{babel}
\usepackage{amsmath}
\usepackage{amsfonts}
\usepackage{amssymb}
\usepackage[margin=2.5cm]{geometry}

\begin{document}

\begin{center}
    \Large \textbf{Evaluación de Examen 2: Unidad II} \\
    \large Física para Ciencias de la Salud (005-1713) \\
    \vspace{0.5cm}
    \textbf{Estudiante:} Bárbara Marcano \quad \textbf{C.I.:} 31.729.255
    \vspace{0.3cm} \\
    \large \textbf{Calificación Final: 6,5/10 puntos}
\end{center}

\vspace{0.5cm}

\section*{Análisis de Resultados}

\subsection*{1. Cinemática (Aceleración media)}
\textbf{Procedimiento y resultado:} La estudiante identifica de forma ordenada los datos iniciales y la ecuación del movimiento ($a = \frac{v_f - v_i}{t}$). La división aritmética $\frac{0,6}{0,15}$ es correcta ($4$), pero comete un error en el análisis dimensional al expresar el resultado con unidades de velocidad ($4 \text{ m/s}$) en lugar de unidades de aceleración ($4 \text{ m/s}^2$). \\
\textbf{Veredicto:} \textbf{Parcialmente correcto} (Penalización por unidades incorrectas) (1,5/2 pts).

\subsection*{2. Dinámica (Leyes de Newton)}
\textbf{Procedimiento y resultado:} Identifica los datos, pero comete un error conceptual grave al relacionar la masa y la fuerza. En lugar de aplicar el despeje analítico correspondiente a la Segunda Ley de Newton ($a = \frac{F}{m}$), plantea erróneamente una multiplicación entre la masa y la fuerza ($25 \text{ kg} \cdot 150 \text{ N}$), arrojando un valor de $3750$. El resultado real de la aceleración era $6 \text{ m/s}^2$. \\
\textbf{Veredicto:} \textbf{Incorrecto} (Se otorgan puntos parciales por la identificación de datos) (0,5/2 pts).

\subsection*{3. Trabajo Mecánico}
\textbf{Procedimiento y resultado:} Plantea correctamente la relación matemática para el cálculo del trabajo mecánico, multiplicando la fuerza ejercida por el desplazamiento ($400 \text{ N} \cdot 0,8 \text{ m}$). El resultado final de $320 \text{ J}$ y su denotación de unidades son perfectamente exactos. \\
\textbf{Veredicto:} \textbf{Correcto} (2/2 pts).

\subsection*{4. Energía Potencial}
\textbf{Procedimiento y resultado:} Extrae de manera estructurada los datos del enunciado y plantea correctamente la fórmula de la energía potencial gravitatoria ($E_p = m \cdot g \cdot h$). Efectúa la multiplicación de los factores (incluyendo $g = 9,8 \text{ m/s}^2$) y obtiene la cifra matemática correcta de $17,64 \text{ J}$. \\
\textbf{Veredicto:} \textbf{Correcto} (2/2 pts).

\subsection*{5. Potencia Mecánica}
\textbf{Procedimiento y resultado:} Existe una confusión conceptual en la definición de la potencia mecánica. La estudiante relaciona erróneamente el trabajo y el tiempo a través de una multiplicación ($60 \text{ s} \cdot 75 \text{ J}$), lo que la lleva a un resultado incorrecto de $4500 \text{ W}$. La ecuación física adecuada dicta que la potencia es el trabajo dividido entre el tiempo ($P = \frac{W}{t}$), por lo que la cifra real era de $1,25 \text{ W}$. \\
\textbf{Veredicto:} \textbf{Incorrecto} (Se otorgan puntos parciales por intentar relacionar las variables del enunciado) (0,5/2 pts).

\vspace{0.5cm}
\section*{Comentarios Generales y Sugerencias}
El examen denota un buen orden inicial al extraer los datos antes de operar (notable en las preguntas 3 y 4), logrando excelentes resultados en el bloque de Energía y Trabajo.
\begin{itemize}
    \item Se aconseja encarecidamente repasar las **fórmulas base de las leyes de Newton y la Potencia**. Las magnitudes deben conservar sus relaciones matemáticas (división frente a multiplicación) para que el resultado tenga sentido físico.
    \item Se debe prestar más atención al **análisis dimensional y a las unidades**. Las unidades siempre deben acompañar el final del problema con exactitud matemática (recordar el cuadrado en los segundos de la aceleración).
\end{itemize}

\end{document}